% Resume-only Projects, in Chinese. A short selection from the CV's much longer
% cv/projects.tex.
%
% Two reasons this section exists on the resume at all. One, with 4D imaging
% radar and the open-source tools removed, the resume had a mostly empty second
% page; this is the most substantive thing available to fill it, and it is
% filled with hardware actually built rather than with awards. Two, none of
% these four items touches either removed topic: they are the doctoral and
% pre-doctoral hardware, which is also the part of the record the Experience
% section cannot show.
%
% The open-source subsection of the CV's Projects (RadarSimX, SensorView,
% Antenna Array Analysis, CommProbe, Tx-Line Calculator, Hexapod) is
% deliberately not carried over here.

\cvsection{代表性项目}

\begin{cventries}

    \cventry
        {便携式 K 波段 MIMO FMCW 雷达}
        {博士研究}
        {}
        {2017}
        {{\cvdescfont 16 发 16 收非均匀间隔阵列，视场约 90\textdegree，方位与俯仰角分辨率约 $\pm$1.5\textdegree，用于三维定位。}}

    \cventry
        {射频波束形成定位雷达}
        {博士研究}
        {}
        {2017}
        {{\cvdescfont 基于 PIN 二极管矢量控制器的 K 波段二维波束形成短距定位雷达，连续扫描范围 $\pm$45\textdegree，整板 11.5\,$\times$\,3.9\,cm、69\,g。}}

    \cventry
        {K 波段多模便携雷达}
        {博士研究}
        {}
        {2015--2017}
        {{\cvdescfont 24\,GHz 便携雷达，可配置为 FMCW 或干涉测量模式，支持定位、距离-多普勒处理、非接触语音感知与风机结构健康监测。}}

    \cventry
        {24\,GHz 雷达收发芯片}
        {与 yearONE LLC 合作}
        {}
        {2013--2014}
        {{\cvdescfont GlobalFoundries 0.13\,\textmu m CMOS 工艺，2\,$\times$\,2\,mm 裸片；负责发射机部分，含 LC 压控振荡器、多相滤波器、射频缓冲器与 QFN 封装匹配网络。}}

\end{cventries}
